\documentclass[11pt]{article}

\usepackage[a4paper,margin=1in]{geometry}
\usepackage[T1]{fontenc}
\usepackage[utf8]{inputenc}
\usepackage{lmodern}
\usepackage{amsmath}
\usepackage{amssymb}
\usepackage{booktabs}
\usepackage{graphicx}
\usepackage{hyperref}
\usepackage{natbib}
\usepackage{setspace}
\usepackage{enumitem}
\usepackage{titlesec}

\hypersetup{
  colorlinks=true,
  linkcolor=blue,
  citecolor=blue,
  urlcolor=blue
}

\setstretch{1.1}
\titleformat{\section}{\large\bfseries}{\thesection.}{0.5em}{}
\titleformat{\subsection}{\normalsize\bfseries}{\thesubsection}{0.5em}{}

\title{\textbf{Hydra-Cool: Buoyancy-Assisted Seawater Retrofit Cooling and Feasible Design Windows for Hyperscale Data Centers}}
\author{Obada Dallo\\
\small Hydra-Cool Open Research Project\\
\small Corresponding author email to be added before submission}
\date{\today}

\begin{document}
\maketitle

\begin{abstract}
Hydra-Cool is a buoyancy-assisted seawater cooling concept intended to reduce the electrical burden of hyperscale data center cooling through deep-water heat rejection, density-driven hydraulic assistance, elevated discharge storage, gravity return flow, and optional micro-hydropower recovery. This study evaluates Hydra-Cool primarily as a retrofit-assist architecture rather than as a universally passive standalone replacement. The simulation framework combines thermal flow requirements, seawater density variation, buoyancy pressure generation, Darcy-Weisbach friction losses, heat-exchanger pressure losses, pump assistance, and optional turbine recovery \citep{fofonoff1983,ebrahimi2014,swamee1976,pal2021,swart2020}. A staged simulation campaign was used to refine the design space through broad screening, candidate pruning, and a focused high-resolution design-window study. The current results show that passive standalone operation is rare, while hybrid retrofit assist dominates the viable solution space. In the focused window, the verified rerun reports a PASS rate of approximately 48.50\%, with 4.39\% of all cases classified as passive standalone and 44.10\% classified as hybrid retrofit assist. The dominant failure mode across the campaign is insufficient velocity, with a smaller secondary contribution from unmet thermal duty. Sensitivity analysis identifies IT load, heat-exchanger pressure drop, number of pipes, pipe diameter, and temperature rise as the principal feasibility drivers. These findings support Hydra-Cool as a constrained retrofit cooling pathway, while also showing that the baseline and auxiliary-load assumptions still require stricter calibration before publication-final savings claims can be defended.
\end{abstract}

\noindent\textbf{Keywords:} data center cooling, seawater cooling, buoyancy-assisted flow, retrofit cooling, hydraulic feasibility

\section{Introduction}

Cooling energy is a major constraint on the growth of hyperscale digital infrastructure, especially as AI-oriented computing density continues to rise \citep{masanet2020,shehabi2024}. Conventional cooling plants rely on mechanically driven chillers, pumps, and cooling towers, all of which impose parasitic electrical demand on the facility \citep{greengrid2011,ebrahimi2014}. Hydra-Cool explores whether part of this burden can be offset through a buoyancy-assisted hydraulic architecture using deep cold seawater, elevated discharge, and gravity return.

The concept is intentionally framed here as a retrofit-assist architecture. The primary question is not whether the loop can replace a conventional plant under fully passive conditions in all cases, but whether it can reduce the net cooling burden when integrated with an existing system.

\section{Related Work}

Hydra-Cool sits at the intersection of several established engineering literature streams. The first is the broad literature on data center cooling systems, including conventional mechanical plants, economizer strategies, and liquid cooling for high-density computing \citep{ebrahimi2014,masanet2020,shehabi2024,khalaj2017}. That literature provides the baseline context against which Hydra-Cool must be judged.

The second stream concerns thermosiphons, natural circulation, and buoyancy-driven heat transport. These systems have been studied for decades because density differences can generate useful circulation under favorable geometric and thermal conditions \citep{zhang2018,swart2020}. This literature is highly relevant to Hydra-Cool, but it is rarely framed in the context of coastal data center retrofit hydraulics.

The third stream is seawater-based cooling and coastal thermal infrastructure. Sea water air conditioning and related marine cooling concepts show that cold seawater can be an effective heat sink where geography and civil design permit \citep{sanjivy2023}. However, those approaches are normally discussed as pumped seawater infrastructure rather than buoyancy-assisted retrofit loops.

Hydra-Cool overlaps with these streams but is not identical to any of them. It is not simply a standard pumped seawater loop, because it seeks hydraulic assistance from buoyancy and gravity. It is not simply a passive thermosiphon, because the dominant viable mode remains \texttt{HYBRID\_RETROFIT\_ASSIST}. The resulting research gap is narrow but meaningful: existing work covers seawater cooling, natural circulation, and thermosiphon behavior, but limited work frames these ideas explicitly as a retrofit-assist hydraulic architecture for large coastal data center cooling systems.

\section{Physical Basis}

The model relies on five standard engineering relationships:

\begin{enumerate}[label=\arabic*)]
  \item Heat-removal requirement:
  \[
  \dot{m} = \frac{P_{\text{thermal}}}{C_p \Delta T}
  \]
  \item Buoyancy pressure:
  \[
  \Delta P_{\text{buoyancy}} = (\rho_{\text{cold}} - \rho_{\text{hot}}) g H
  \]
  \item Hydraulic losses:
  \[
  \Delta P_{\text{losses}} = \Delta P_{\text{friction}} + \Delta P_{\text{minor}} + \Delta P_{\text{HX}}
  \]
  \item Hybrid total power:
  \[
  P_{\text{Hydra,total}} = P_{\text{pump,assist}} + P_{\text{aux}} - P_{\text{turbine}}
  \]
  \item Retrofit savings fraction:
  \[
  \eta_{\text{savings}} = \frac{P_{\text{baseline}} - P_{\text{Hydra,total}}}{P_{\text{baseline}}}
  \]
\end{enumerate}

Seawater properties are approximated using UNESCO algorithms \citep{fofonoff1983}. Hydraulic resistance is modeled through Darcy-Weisbach with an explicit turbulent friction-factor approximation \citep{swamee1976}. Water-hammer interpretation follows contemporary numerical review literature \citep{pal2021}, while natural-circulation interpretation is informed by closed-loop thermosyphon analysis \citep{swart2020,zhang2018}.

\section{Methods}

\subsection{Operating layers}

Three operating layers are tracked explicitly:

\begin{itemize}
  \item \texttt{PASSIVE\_STANDALONE}
  \item \texttt{HYBRID\_RETROFIT\_ASSIST}
  \item \texttt{TURBINE\_RECOVERY\_ACTIVE}
\end{itemize}

The turbine layer is treated as optional and secondary to the retrofit cooling objective.

\subsection{Stage-based simulation campaign}

The campaign was organized into three stages:

\begin{table}[h]
\centering
\begin{tabular}{lll}
\toprule
Stage & Purpose & Scenario count \\
\midrule
1 & Broad screening & 24,000 \\
2 & Candidate pruning & 543 \\
3 & Focused design window & 995,328 \\
\bottomrule
\end{tabular}
\caption{Stage structure of the Hydra-Cool simulation campaign.}
\end{table}

The focused window emphasized pipe diameters near 0.8--1.5 m, high vertical lift, elevated $\Delta T$, low-to-moderate heat-exchanger loss, and IT load in the 50--200 MW range.

\subsection{Classification}

Scenarios are classified as:

\begin{itemize}
  \item \textbf{PASS}: hydraulically feasible, satisfies the thermal duty, and achieves at least 10\% retrofit energy reduction
  \item \textbf{MARGINAL}: hydraulically feasible, satisfies the thermal duty, and achieves at least 2\% but less than 10\% reduction
  \item \textbf{FAIL}: hydraulically infeasible, unable to sustain the thermal duty, or insufficient in energy benefit
\end{itemize}

Scenarios fail explicitly when useful minimum velocity cannot be sustained or when achievable mass flow cannot support the required thermal duty. When both are present, insufficient velocity is reported as the primary failure mode.

\section{Results}

\subsection{Stage 1}

Stage 1 broad screening produced a PASS rate of 8.65\% and a FAIL rate of 91.35\%. Passive standalone cases did not constitute a meaningful success mode in this broad phase.

\subsection{Stage 2}

Stage 2 candidate pruning reduced the active development space to 543 scenarios and raised the PASS rate to 24.13\%. The qualitative result remained unchanged: hybrid retrofit assist dominated successful outcomes.

\subsection{Stage 3}

The focused high-resolution window yielded the clearest current feasibility signal:

\begin{itemize}
  \item PASS rate: 48.50\%
  \item FAIL rate: 51.50\%
  \item Passive standalone PASS rate: 4.39\%
  \item Hybrid retrofit-assist PASS rate: 44.10\%
\end{itemize}

The strongest-performing cases clustered around high IT load, two-pipe configurations, pipe diameters near 1.0--1.2 m, high vertical lift, high $\Delta T$, and low heat-exchanger loss.

\subsection{Failure modes and sensitivity}

Across the campaign, insufficient velocity was the dominant hydraulic failure mechanism. Sensitivity analysis in the focused window ranked the principal variables as:

\begin{enumerate}[label=\arabic*)]
  \item IT load
  \item Heat-exchanger pressure drop
  \item Number of pipes
  \item Pipe diameter
  \item $\Delta T$
\end{enumerate}

\section{Benchmark Positioning}

To strengthen research framing, Hydra-Cool was also positioned against a conventional mechanically driven cooling plant, a pumped seawater loop without buoyancy assistance, and uncommon passive-natural edge cases. This benchmark is intentionally conservative and normalized. It should be read as a positioning device rather than a calibrated techno-economic comparison.

\begin{table}[h]
\centering
\begin{tabular}{lllll}
\toprule
System & Relative burden & Pump dependence & Passive contribution & Maturity \\
\midrule
Conventional plant & 1.00 & Very high & Very low & High \\
Pumped seawater loop & 0.70 & High & Low & Moderate to high \\
Hydra-Cool hybrid assist & 0.30 & Moderate & High & Low / research-stage \\
Hydra-Cool passive edge case & 0.18 & Very low & Very high & Very low \\
\bottomrule
\end{tabular}
\caption{Normalized benchmark framing used to position Hydra-Cool relative to known alternatives.}
\end{table}

\section{Limitations}

The current work remains a research-stage feasibility study. Major limitations include:

\begin{itemize}
  \item simplified baseline cooling assumptions,
  \item lumped auxiliary loads,
  \item steady-state rather than transient operation,
  \item absence of site-specific marine civil constraints,
  \item lack of external calibration against operating hyperscale cooling plants.
\end{itemize}

Accordingly, the present savings estimates should be interpreted as provisional rather than publication-final.

\section{Discussion}

The present results support Hydra-Cool primarily as a hybrid retrofit-assist concept. The staged workflow indicates that the design window is non-zero, but also that its feasibility only becomes clear after the parameter space is progressively narrowed. This means the research should be interpreted as evidence for a constrained but real design window, not as proof of a universally passive cooling architecture.

Insufficient velocity dominates failure because buoyancy alone does not guarantee useful circulation. The most important parameters all influence whether buoyancy can be converted into sustained thermal transport rather than dissipated through hydraulic losses. This explains why heat-exchanger pressure drop, pipe count, and pipe diameter matter so strongly in the focused design window.

The hybrid retrofit-assist framing is therefore the most credible interpretation. A cooling retrofit does not need to eliminate pumping entirely to be valuable; it only needs to reduce the net burden of active heat rejection in a repeatable and defensible way. In that sense, the benchmark comparison strengthens context but does not replace validation. Hydra-Cool should currently be understood as a simulation-based research concept whose claims remain limited by baseline uncertainty, steady-state modeling, and the absence of prototype-scale evidence.

\section{Conclusion}

Hydra-Cool has a non-zero feasible design window and is most credible as a buoyancy-assisted retrofit cooling architecture. The dominant success mode is hybrid retrofit assist, while passive standalone operation is rare. The principal engineering bottleneck is not the existence of buoyancy, but the ability to sustain useful circulation velocity under realistic hydraulic losses. The next phase of the work should emphasize stricter baseline calibration rather than wider random scenario generation.

\section*{Data and Figure Availability}

The primary simulation summaries and publication figures are stored in the repository output directory. A Markdown figure board is available in \texttt{docs/PUBLICATION\_FIGURES.md}.

\bibliographystyle{unsrtnat}
\bibliography{references}

\end{document}
